\documentclass[a4paper,10pt]{article}
\usepackage[utf8]{inputenc}
\usepackage{graphicx}
\usepackage{subfigure}
\usepackage{fancyhdr}
\usepackage{listings}
\usepackage{ gensymb }


%opening

\pagestyle{fancy} %eigener Seitenstil
\fancyhf{} %alle Kopf- und Fußzeilenfelder bereinigen
\fancyhead[L]{SEP Abgabegespräch SS2015/16} %Kopfzeile links
\fancyhead[C]{} %zentrierte Kopfzeile
\fancyhead[R]{\includegraphics[width=1cm]{../../template/tugraz_bw}} %Kopfzeile rechts
\renewcommand{\headrulewidth}{0.4pt} %obere Trennlinie
\fancyfoot[C]{} %Seitennummer
\renewcommand{\footrulewidth}{0.4pt} %untere Trennlinie


\lstset{language=C++, basicstyle=\ttfamily\tiny}  

\begin{document}

\section*{Mikrobensimulator}

Schreibe ein Programm, welches das Leben von Mikroben (Bakterien und Viren) simuliert. Dazu soll eine Liste von Mikroben aus einer Datei eingelesen werden. Jede Zeile hat das Format \texttt{<typ> <name> <anzahl> <faktor>}. Der Typ ist entweder \texttt{B} für Baktieren oder \texttt{V} für Viren. Der Name ist ein beliebiger String, die Anzahl gibt an wie viele Mikroben es von dieser Sorte am Anfang gibt. Der Faktor bestimmt das Wachstum der Mikroben. In jeder Zeiteinheit vermehren sich die Mikroben:
\begin{description}
 \item[Bakterien] Die Anzahl der Bakterien wächst um den angegebenen Faktor, d.h. er wird zur aktuellen Anzahl addiert
 \item[Viren] Die Anzahl der Viren vervielfacht sich um den angegebenen Faktor, d.h. er wird mit der aktuellen Anzahl multipliziert
\end{description}

Das Programm soll so viele Zeiteinheiten wie beim Aufruf angegeben simulieren. Dann soll von jedem Bakteriums- und Virentyp ausgegeben werden, wie viele es nach dieser Zeiteinheit gibt (siehe Beispiele). Bei der Ausgabe sollen Kommazahlen immer auf die nächste Ganzzahl abgerundet ausgegeben werden.

Ist die Anzahl der Parameter falsch, soll \\
``\texttt{Usage: ./microbes <microbes~list> <time>}'' ausgegeben werden.

\vspace{1em}

Beispiel für eine Mikroben-Liste \texttt{microbes.txt}:
\begin{lstlisting}[frame=single]
V Ebola 1 3
V Smallpox 2 1.7
V Influenza 100 1.1
B Lactobacillus 10 25
B Salmonella 1000 1.8
B Streptococcus 0 100
\end{lstlisting}

\subsection*{Beispiel}
\footnotesize
\begin{verbatim}
./microbes microbes.txt 0
Ebola (1)
Smallpox (2)
Influenza (100)
Lactobacillus (10)
Salmonella (1000)
Streptococcus (0)

./microbes microbes.txt 10
Ebola (59049)
Smallpox (403)
Influenza (259)
Lactobacillus (260)
Salmonella (1017)
Streptococcus (1000)
\end{verbatim}

\subsection*{Aufruf}
\texttt{./microbes <microbes list> <time>}

\subsection*{Bewertung}
\begin{itemize}
 \item Dokumentation und Programmierstil: 1 Punkt
 \item Einhalten der OOP Konzepte: 1 Punkt
 \item Struktur \& Korrektheit: 1 Punkt
\end{itemize}

\newpage


\end{document}
